\documentclass[spanish,letterpaper,oneside]{article}

\def\documenttitle {¿Qué sabemos y cómo lo enseñamos? · Ficha de trabajo}
\def\documentsubtitle {Temas selectos de matemáticas I · Cohorte SD26}
\def\documentsubject {Actividad 1 · Material para estudiantes}
\def\documentauthor {Mtro. Martín Jonathan de la Cruz Muñoz}
\def\coursename {Temas selectos de matemáticas I}
\def\coursecode {PFM02}
\def\universityname {Instituto de Investigación, Innovación y Estudios de Posgrado para la Educación del Estado de Nuevo León}
\def\universityfaculty {Maestría en Aprendizaje y Enseñanza de las Matemáticas}
\def\universitydepartment {Temas selectos de matemáticas I}
\def\universitydepartmentimage {img/iiiepe/marca-iiiepe.png}
\def\universitydepartmentimagecfg {height=1.75cm}
\def\universitylocation {Monterrey, Nuevo León}
\def\studentfullname {Mtro. Martín Jonathan de la Cruz Muñoz}
\def\studentid {00874}
\def\studentemail {martin.delacruz@campus.iiiepe.edu.mx}
\def\programperiod {2do periodo, agosto--noviembre de 2026 (SD26)}
\def\programmodality {Modalidad presencial y en línea}
\def\programshort {MAEM}

\def\authortable {
  \begin{tabular}{ll}
    \textbf{Docente:} & \studentfullname \\
    Matrícula: & \studentid \\
    Unidad: & \coursename \\
    Periodo: & \programperiod \\
    \multicolumn{2}{l}{Fecha: 10 de octubre de 2026}
  \end{tabular}
}

\input{template}

\usepackage{helvet}
\renewcommand{\familydefault}{\sfdefault}
\usepackage{amsmath,amssymb,array,longtable,tikz,tabularx,colortbl,enumitem}
\definecolor{iiiepeNavy}{HTML}{163B5C}
\definecolor{iiiepeBlue}{HTML}{1F6F9F}
\definecolor{iiiepePale}{HTML}{EEF5F8}
\def\portraittitlecolor {iiiepeNavy}

\setlength{\parskip}{5pt}\setlength{\parindent}{0pt}
\begin{document}\templatePortrait\templatePagecfg\templateFinalcfg\fancyfoot{}\fancyfoot[L]{\scriptsize IIIEPE · MAEM · SD26}\fancyfoot[C]{\scriptsize\thepage}\fancyfoot[R]{\scriptsize Ficha Act. 1}\setstretch{1.2}\raggedbottom

\section*{Explorar áreas y perímetros desde el desafío}
Nombre: \rule{8cm}{.2pt}\hfill Fecha: \rule{2.5cm}{.2pt}

Resuelve cada desafío registrando tus tablas, dibujos, comprobaciones y explicaciones matemáticas.

\subsection*{Reto 1. ¡Construye el jardín más grande!}
La escuela dispone de exactamente $24\text{ m}$ de cerca para construir un jardín rectangular de dimensiones enteras.
\begin{enumerate}
  \item Completa todas las combinaciones distintas $(l, a)$ que cumplen $2(l+a)=24$ y calcula su área $A=l\cdot a$.
  \item ¿Cuántos diseños diferentes encontraste? ¿Cuál logra la mayor área? Explica cómo sabes que no existe otro rectángulo entero con mayor superficie.
\end{enumerate}
\vspace{2.0cm}

\subsection*{Reto 2. El arquitecto bajo presión}
Se solicita un jardín rectangular de $48\text{ m}^2$ con solo $28\text{ m}$ de cerca. ¿Es posible construirlo? Demuestra tus dimensiones posibles y argumenta por qué convencerías al director con razones matemáticas.
\vspace{2.0cm}

\subsection*{Reto 3. El jardín junto a una pared}
Si una pared existente cubre uno de los lados largos del jardín y solo cercas los otros tres lados con $24\text{ m}$ de malla ($2a+l=24$), ¿es posible alcanzar al menos $70\text{ m}^2$ con dimensiones enteras? Indica todas las opciones válidas y cuál maximiza el área.
\vspace{2.0cm}

\end{document}
